% META: {"id": "TCOMPL-ATT-EX-037", "chapitre": "TCOMPL-TEMPS-ATTENTE", "type_objet": "exercice", "capacites": ["TCOMPL-ATT-C7"], "capacites_codes": ["C7"], "methodes": [], "parcours": 1, "competences": ["calculer"], "duree_min": 15, "mode_creation": "ex_nihilo", "sources_inspiration": [], "fichier_tex": "chapitres/TCOMPL-TEMPS-ATTENTE/exercices/TCOMPL-ATT-EX-037.tex", "corrige_tex": "chapitres/TCOMPL-TEMPS-ATTENTE/corriges/TCOMPL-ATT-CO-037.tex", "status": "needs_review"}
% BEGIN-VERIFY
% from sympy import symbols, integrate, Rational, simplify
% x = symbols('x', real=True)
% a, b = 0, 15
% f = Rational(1, b - a)
% assert integrate(f, (x, a, b)) == 1
% assert integrate(x * f, (x, a, b)) == Rational(15, 2)
% moment2 = integrate(x**2 * f, (x, a, b))
% assert simplify(moment2 - Rational(15, 2)**2) == Rational(75, 4)
% assert simplify(Rational((b - a)**2, 12) - Rational(75, 4)) == 0
% assert integrate(f, (x, 0, 5)) == Rational(1, 3)
% assert integrate(f, (x, 10, 15)) == Rational(1, 3)
% END-VERIFY
\begin{exercice}{TCOMPL-ATT-EX-037}{1}{15}
Un métro passe toutes les $15$ minutes. Un voyageur arrive à la station sans connaître l'horaire : son temps d'attente $X$, en minutes, suit la loi uniforme sur $[0\,;15]$.
\begin{enumerate}
  \item Donner la densité de $X$ et vérifier qu'il s'agit bien d'une densité de probabilité.
  \item Déterminer la fonction de répartition $F$ de $X$ sur $[0\,;15]$.
  \item Calculer $P(X \leq 5)$ puis $P(10 \leq X \leq 15)$. Que remarque-t-on ?
  \item Calculer $E(X)$ et $V(X)$, et interpréter $E(X)$.
\end{enumerate}
\end{exercice}
